\section*{الفصل \ref{ch:b2:mass-spec-atomic} --- مطيافية الكتلة والمطيافية الذرية}

\begin{solution}{exo:b2:mass-spec-atomic:1}
78؛ $400/2 = 200$؛ $194 + 1 = 195$، أي $[\mathrm M + \mathrm H]^+$.
\end{solution}

\begin{solution}{exo:b2:mass-spec-atomic:2}
كتلة فردية: عدد فردي من ذرات النيتروجين، وهو ذرة واحدة لجزيء بهذا الصغر. \ce{C3H7NO} (البروبانأميد، 73)
أو \ce{C4H11N} (البيوتان-1-أمين، 73).
\end{solution}

\begin{solution}{exo:b2:mass-spec-atomic:3}
متقاربتان: ذرة بروم واحدة. 3 : 1: ذرة كلور واحدة.
\end{solution}

\begin{solution}{exo:b2:mass-spec-atomic:4}
الأصفر: الصوديوم، \qty{589.0}{nm} و\qty{589.6}{nm}. والأحمر القرمزي: الليثيوم، \qty{670.8}{nm}.
% ledger: asd:NaD2, asd:NaD1, asd:Li671
\end{solution}

\begin{solution}{exo:b2:mass-spec-atomic:5}
$n_C \approx 6.7/1.11 \approx 6$.
\end{solution}

\begin{solution}{exo:b2:mass-spec-atomic:6}
\ce{CO}: 27.9949؛ \ce{N2}: 28.0061؛ \ce{C2H4}: 28.0313 u. قدرات الميز: $\ce{CO}/\ce{N2}$ $28/0.0112
\approx 2.5 \times 10^3$؛ $\ce{N2}/\ce{C2H4}$ $28/0.0252 \approx 1.1 \times 10^3$؛ $\ce{CO}/\ce{C2H4}$
$28/0.0364 \approx 7.7 \times 10^2$.
% ledger: mass:C12, mass:O16, mass:N14, mass:H1
\end{solution}

\begin{solution}{exo:b2:mass-spec-atomic:7}
86: \omterm{def:b2:mass-spec-atomic:molecular-ion}{الأيون الجزيئي} للمركب \ce{C5H10O}. 57: \ce{C2H5CO+}، أيون أسيليوم ناتج من الانشطار $\alpha$ (فقد
جذر إيثيل، 29). 29: \ce{C2H5+}.
% ledger: ms:pentan-3-one
\end{solution}

\begin{solution}{exo:b2:mass-spec-atomic:8}
86: $\mathrm M^{+\bullet}$. 71: $\mathrm M - \ce{CH3}$، \omterm{def:b2:aromatic-substitution:friedel-crafts}{أيون الأسيليوم} \ce{C3H7CO+}. 43: \ce{CH3CO+}
(فقد جذر البروبيل). 58: \omterm{def:b2:mass-spec-atomic:fragmentation}{إعادة ترتيب ماكلافرتي}، عبر هيدروجين على الكربون $\gamma$
في سلسلة البروبيل، مع فقد الإيثين. وفي البنتان-3-ون ليس لمجموعتي الإيثيل كربون $\gamma$: فلا
أيون ماكلافرتي؛ وقمته الصغيرة عند 58 هي قمة النظير \ce{^{13}C} للأيون عند 57 (ثلاث ذرات كربون،
نحو 3.3~\%).
% ledger: ms:pentan-2-one, ms:pentan-3-one
\end{solution}

\begin{solution}{exo:b2:mass-spec-atomic:9}
خط مار بالمبدأ: الميل $\sum c_iA_i/\sum c_i^2 = 1.553/30.0 \approx \qty{0.0518}{L/mg}$. والعينة:
$0.128/0.0518 \approx \qty{2.47}{mg/L}$ من النحاس.
\end{solution}

\begin{solution}{exo:b2:mass-spec-atomic:10}
$t = L\sqrt{m/(2eU)} = 1.00\sqrt{1000 \times 1.6605 \times 10^{-27}/(2 \times 1.6022 \times 10^{-19}
\times 2.00 \times 10^4)} \approx \qty{16.1}{\micro s}$. ومن أجل 1001 u، يزداد $t$ بالعامل
$\sqrt{1.001} \approx 1 + 0.0005$: $\Delta t \approx \qty{8.0}{ns}$.
% ledger: const:u, const:e
\end{solution}

\begin{solution}{exo:b2:mass-spec-atomic:11}
ذرتا كلور: $(0.758 + 0.242x)^2$ يعطي $0.575$، $0.367$، $0.0586$: قيم $m/z$ 84 و86 و88 بالنسبة
100 : 64 : 10.
\end{solution}

\begin{solution}{exo:b2:mass-spec-atomic:12}
\ce{^{40}Ar^{16}O+}: $39.9624 + 15.9949 = 55.9573$؛ \ce{^{56}Fe}: 55.9349؛ $R = 56/0.0224 \approx 2.5 \times
10^3$. وإلا فقِس نظيرًا آخر للحديد، \ce{^{57}Fe}، أو أزل \ce{ArO+} في خلية تصادم
قبل المحلل.
% ledger: mass:Ar40, mass:Fe56, mass:O16
\end{solution}

\begin{solution}{pb:b2:mass-spec-atomic:1}
\textbf{1.} 156 و158 و160 (مع 157 و159): يوجد الجزيء في عدة صور نظائرية.
\textbf{2.} ثلاث قمم تفصل بينها وحدتان، والوسطى أكبرها: ذرة كلور واحدة وذرة بروم واحدة.
\textbf{3.} 156.
\textbf{4.} كتلة زوجية: لا نيتروجين، أو عدد زوجي منه.
\textbf{5.} $156 - 79 - 35 = 42$: \ce{C3H6}.
\textbf{6.} \ce{C3H6BrCl}؛ ودرجة عدم التشبع $(2 \times 3 + 2 - 6 - 2)/2 = 0$: لا حلقة، ولا رابطة
مزدوجة.
\textbf{7.} $0.5/14.5 \approx 3.4~\%$، $3.4/1.11 \approx 3$ ذرات كربون.
\textbf{8.} القمة عند 157 معطاة بدقة 0.1 على سلم قيمتها فيه 0.5: أي ارتياب قدره 20~\%.
\textbf{9.} \ce{C3H6^{81}Br^{35}Cl} و\ce{C3H6^{79}Br^{37}Cl}.
\textbf{10.} $3 \times 12 + 6 \times 1.007825 + 78.918338 + 34.968853 \approx 155.9341$.
\textbf{11.} $156/(156.151 - 155.934) \approx 7.2 \times 10^2$.
\textbf{12.} نعم: بهيدروجينين أقل، 154.
\textbf{13.} فُقدت ذرة بروم (79 أو 81): \ce{C3H6Cl+}، عند 77 مع \ce{^{35}Cl}، وعند 79 مع \ce{^{37}Cl}،
3 : 1.
\textbf{14.} فقد HBr: الكاتيون الجذري $\mathrm{C_3H_5Cl^{+\bullet}}$، زوجي الكتلة.
\textbf{15.} \ce{C3H5+}، كاتيون الأليل، بعد فقد الهالوجينين كليهما (Br، ثم HCl).
\textbf{16.} \ce{CH2Cl+} (49 و51، 3 : 1).
\textbf{17.} \ce{CH2Br+} (93 و95) و\ce{C2H4Br+} (107 و109)، كل منهما 1 : 1: البروم على أحد طرفي
السلسلة، والثاني ناتج من فقد \ce{CH2Cl}.
\textbf{18.} الرابطة \ce{C-Br} أضعف من الرابطة \ce{C-Cl}، و\ce{Br} الجذر المغادر
الأفضل.
\textbf{19.} لا: لا قمم عند 127--131؛ و\ce{CH2Cl+} و\ce{CH2Br+} يبيّنان أن كل هالوجين على \ce{CH2}
طرفية.
\textbf{20.} \ce{BrCH2CH2CH2Cl}، أي 1-برومو-3-كلوروبروبان.
\textbf{21.} 156/158/160: M؛ 77/79: M $-$ Br؛ 76: M $-$ HBr؛ 107/109: M $-$ \ce{CH2Cl}؛
93/95: \ce{CH2Br+}؛ 49/51: \ce{CH2Cl+}؛ 41: \ce{C3H5+}؛ 39، 27: \ce{C3H3+}، \ce{C2H3+}.
\textbf{22.} لا: فلا مجموعة كربونيل فيه.
\textbf{23.} 1-برومو-2-كلوروبروبان، \ce{CH3CHClCH2Br}: فقد \ce{CH2Br} يعطي \ce{CH3CHCl+} عند 63/65.
\textbf{24.} \ce{C3H6^{81}Br^{37}Cl}.
\textbf{25.} $M$: $0.758 \times 0.5069 = 0.384$؛ $M+2$: $0.758 \times 0.4931 + 0.242 \times 0.5069 =
0.496$؛ $M+4$: $0.242 \times 0.4931 = 0.119$. النسبة $\boldsymbol{\approx 3.2 : 4.2 : 1}$، أي نحو 3 : 4 :
1؛ ويعطي الطيف $14.5 : 18.6 : 4.4 = 3.3 : 4.2 : 1$.
\end{solution}
